% Part: model-theory
% Chapter: models-of-arithmetic
% Section: introduction

\documentclass[../../../include/open-logic-section]{subfiles}

\begin{document}

% SEGMENT OLP-0192-S01

\olfileid{mod}{mar}{int}
\section{அறிமுகம்}

எண்கணிதத்தின் \emph{நிலையான மாதிரி} என்பது
$\Domain{N} = \Nat$ ஆக உள்ள !!{structure}~$\Struct{N}$; அதில்
$\Obj{0}$, $\prime$, $+$, $\times$, $<$ ஆகியவை எதிர்பார்க்கும்
வகையில் பொருள்கொள்ளப்படுகின்றன. அதாவது, $\Obj{0}$ என்பது $0$;
$\prime$ என்பது அடுத்தெண் சார்பு; $+$ என்பது~$\Nat$ இலுள்ள எண்களின்
கூட்டலாகவும் $\times$ அவற்றின் பெருக்கலாகவும் பொருள்கொள்ளப்படுகின்றன.
குறிப்பாக,
\begin{align*}
  \Assign{\Obj{0}}{N} & = 0\\
  \Assign{\prime}{N}(n) & = n + 1\\
  \Assign{+}{N}(n, m) & = n + m\\
  \Assign{\times}{N}(n, m) & = nm
\end{align*}
ஆகும். உறுதியாக,~$\Nat$ அல்லாத ஆட்களங்களைக் கொண்ட
$\Lang{L_A}$-கட்டமைப்புகளும் உள்ளன. எடுத்துக்காட்டாக,
$\Domain{M} = \{a\}^*$ (ஒரே குறியீடு~$a$ இன் இறுதிக்குட்பட்ட
தொடர்கள், அதாவது $\emptyset$, $a$, $aa$, $aaa$, \dots) என்ற
ஆட்களமுடைய $\Struct{M}$ ஐயும் பின்வரும் பொருள்கோள்களையும் கொள்ளலாம்:
\begin{align*}
  \Assign{\Obj{0}}{M} & = \emptyset\\
  \Assign{\prime}{M}(s) & = s \concat a\\
  \Assign{+}{M}(n, m) & = a^{n + m}\\
  \Assign{\times}{M}(n, m) & = a^{nm}
\end{align*}
இவ்விரு கட்டமைப்புகளும் ``சாரமாக ஒன்றே'': !!{domain}s இன்
!!{element}s மட்டுமே வேறுபடுகின்றன; பொருள்கோள் சார்புகளால் அந்த
!!{domain}s இன் !!{element}s ஒன்றுக்கொன்று தொடர்புபடுத்தப்படும் முறை
வேறுபடுவதில்லை. இவ்விரு !!{structure}s
\emph{ஓரினமானவை} என்று கூறுகிறோம்.

~$\Struct{N}$ இல் மெய்யான எந்தக் கோட்பாட்டிற்கும்~$\Struct{N}$ க்கு
ஓரினமல்லாத மாதிரிகளும் உள்ளன என்பது இறுதிக்குட்பட்ட தன்மைத்
தேற்றத்தின் எளிய விளைவு. அத்தகைய கட்டமைப்புகள் \emph{நிலையற்றவை}
எனப்படுகின்றன. ஒரு நிலையான மாதிரியின் (அதாவது $\Struct{N}$ இன்,
ஆனால் அதற்கு ஓரினமான எல்லா !!{structure}s இன் கூட) !!{element}s
நிலையான எண்குறிகள்~$\num{n}$ இன் மதிப்புகளால் முழுவதும் அடங்குகின்றன;
அதாவது,
\[
\Domain{N} = \Setabs{\Value{\num{n}}{N}}{n \in \Nat}
\]
என்பது அவற்றின் சுவையான தன்மை. நிலையற்ற மாதிரிகளில் அவ்வாறில்லை:
$\Struct{M}$ நிலையற்றது எனில், எல்லா~$n$ க்கும்
$x \neq \Value{\num{n}}{M}$ ஆகுமாறு குறைந்தபட்சம் ஒரு
$x \in \Domain{M}$ உள்ளது.

இந்நிலையற்ற உறுப்புகள் மிகவும் சுவையானவை: அவை ``முடிவிலி இயல்
எண்கள்''. ஆனால் அவற்றின் இருப்பு, ஒரு பொருளில், நிறைவின்மை நிகழ்வுகளையும்
விளக்குகிறது. எடுத்துக்காட்டாக, பியானோ எண்கணிதத்தின் முரணின்மைக் கூற்று
$\OCon[\Th{PA}]$, அதாவது $\lnot
\lexists[x][\OPrf[\Th{PA}](x, \gn{\lfalse})]$ ஐக் கருதுக.
$\Th{PA}$ ஆனது $\OCon[\Th{PA}]$ ஐயும் $\lnot \OCon[\Th{PA}]$ ஐயும்
நிரூபிக்காததால், அவற்றுள் எதையும் $\Th{PA}$ உடன் முரணின்றிச் சேர்க்கலாம்.
$\Th{PA}$ முரணின்மையுடையது என்பதால், $\Sat{N}{\OCon[\Th{PA}]}$;
அதன் விளைவாக $\Sat/{N}{\lnot
  \OCon[\Th{PA}]}$. எனவே $\Struct{N}$ ஆனது $\Th{PA}
\cup \{\lnot \OCon[\Th{PA}]\}$ இன் மாதிரி \emph{அல்ல}; அதன் எல்லா
மாதிரிகளும் நிலையற்றவையாக இருக்க வேண்டும். $\Th{PA} \cup
\{\lnot \OCon[\Th{PA}]\}$ இன் மாதிரிகள்,
$\lexists[x][\OPrf[\Th{PA}](\gn{\lfalse})]$ ஐ மெய்யாக்கும் சாட்சியாகப்
பணியாற்றும் ஓர் !!{element} ஐக் கொண்டிருக்க வேண்டும்; அதாவது
~$\Th{PA}$ இலிருந்து ஒரு முரண்பாட்டிற்கான !!a{derivation} இன்
G\"odel எண். அத்தகைய !!a{element} நிலையானதாக இருக்க முடியாது---ஏனெனில்
ஒவ்வொரு~$n$ க்கும் $\Th{PA} \Proves \lnot
\OPrf[\Th{PA}](\num{n}, \gn{\lfalse})$.

\end{document}
